\input{sol_pre.tex}
\input{q_2025-12-30.tex}
\begin{document}
\begin{center}
{\color{msblue}\LARGE\bfseries MILLIY SERTIFIKAT ORIGINAL SAVOLLAR}\\[3pt]
{\color{msblue}\large\bfseries \curdate \qquad Matematika \qquad Yeçimlar}\\[4pt]
{\small Çizmalarda \textcolor{aux}{\bfseries qizil rangdagi} çiziqlar — yeçim uçun qöşimça yasaşlar.}
\end{center}

%================= 1
\Savol{1}
$\underbrace{22\ldots2}_{n}\cdot11=\underbrace{22\ldots2}_{n}0+\underbrace{22\ldots2}_{n}$: çetki raqamlar $2$, oraliqdagi $n-1$ ta raqam $2+2=4$ (masalan, $222\cdot11=2442$). $n=50$ da: $2\underbrace{44\ldots4}_{49}2$.
\Javob{C}

%================= 2
\Savol{2}
$\frac13+\frac14=\frac7{12}$, \ $\frac13-\frac14=\frac1{12}$:
\[\frac{4-\frac{12}7}{4+12}=\frac{16/7}{16}=\frac17.\]
\Javob{B}

%================= 3
\Savol{3}
Jami: $300+315+330+345+390=1680$. Sotilgandan keyin qovun $=3\cdot$tarvuz, ya'ni qolgan mahsulot $4$ ga bölinadi. $1680$ $4$ ga bölinadi, demak sotilgan maşinadagi son ham $4$ ga bölinadi: faqat $300$.\\
Qolgan $1380$: tarvuz $\frac{1380}4=345$ (törtinçi maşina), qovun $315+330+390=1035$.
\Javob{B}

%================= 4
\Savol{4}
Quruq modda: $42\cdot0{,}08=3{,}36$ kg. Qoqida u $84\%$: \ $3{,}36:0{,}84=4$ kg.
\Javob{C}

%================= 5
\Savol{5}
\[\frac{2^{\frac14}\cdot2^{\frac13}\cdot3^{\frac14}}{2^{2}\cdot3^{\frac16}\cdot3^{\frac12}}=2^{\frac14+\frac13-2}\cdot3^{\frac14-\frac16-\frac12}=2^{-\frac{17}{12}}\cdot3^{-\frac5{12}}.\]
$n$-darajada ratsional bölişi uçun $\frac{17n}{12}$ va $\frac{5n}{12}$ butun bölişi kerak $\Rightarrow$ $n=12$.
\Javob{A}

%================= 6
\Savol{6}
$a=\sqrt[3]{\sqrt5+2}$, $b=\sqrt[3]{\sqrt5-2}$: \ $ab=\sqrt[3]{5-4}=1$, \ $a^3-b^3=4$. $x=a-b$ desak,
\[x^3=a^3-b^3-3ab(a-b)=4-3x\ \Rightarrow\ x^3+3x-4=0\ \Rightarrow\ (x-1)(x^2+x+4)=0\ \Rightarrow\ x=1.\]
\Javob{A}

%================= 7
\Savol{7}
$52\pm6\sqrt{43}=(\sqrt{43}\pm3)^2$, şuning uçun ifoda
\[(\sqrt{43}+3)^3-(\sqrt{43}-3)^3=2\left(3\cdot43\cdot3+27\right)=2\cdot414=828.\]
\Javob{D}

%================= 8
\Savol{8}
$ab-c>0$ $\Rightarrow$ $1$; \ $ac-b<0$ $\Rightarrow$ $-1$; \ $bc-a<0$ $\Rightarrow$ $-1$. \ Yiğindi: $1-1-1=-1$.
\Javob{D}

%================= 9
\Savol{9}
$S_8$, $S_7$ — $t^2-26t+144=0$ ning ildizlari: $18$ va $8$. Ösuvçi progressiyada $a_8=S_8-S_7$; $S_8=18$, $S_7=8$ bölsa $a_8=10$ va $a_4=\frac{S_7}7=\frac87<a_8$ — mos (aksi $a_8=-10<a_4$ böladi).
\[S_{15}=15a_8=150.\]
\Javob{A}

%================= 10
\Savol{10}
$y=xq$, $z=xq^2$; \ $2\cdot2y=x+3z$ $\Rightarrow$ $4q=1+3q^2$ $\Rightarrow$ $q=1$ yoki $q=\frac13$. Kamayuvçi: $q=\frac13$.
\Javob{D}

%================= 11
\Savol{11}
$t=x-2$: surat $t^4+(t+1)^2-1=t^4+t^2+2t=t(t^3+t+2)=t(t+1)(t^2-t+2)$,\\
maxraj $t(t-1)+2=t^2-t+2$. \ Natija: $t(t+1)=(x-2)(x-1)=x^2-3x+2$.
\Javob{D}

%================= 12
\Savol{12}
$x>0$: $x^2+10x-11=0$ $\Rightarrow$ $x=1$. \quad $x<0$: $-x^2+10x+11=0$ $\Rightarrow$ $x=-1$ ($x=11$ mos emas).\\
$2$ ta butun yeçim.
\Javob{A}

%================= 13
\Savol{13}
Viyet: $x_1\cdot12x_1=3$ $\Rightarrow$ $x_1=\pm\frac12$; \ $13x_1=\frac{10-a}2$ $\Rightarrow$ $a=10-26x_1=-3$ yoki $23$. Yiğindi $20$.
\Javob{A}

%================= 14
\Savol{14}
$8+\dfrac8{2^x}=x+x\cdot2^x$ $\Rightarrow$ $\dfrac{8(2^x+1)}{2^x}=x(2^x+1)$ $\Rightarrow$ $x\cdot2^x=8$.\\
$x\le0$ da çap tomon manfiy yoki $0$, $x>0$ da ösuvçi $\Rightarrow$ yagona ildiz $x=2$.
\Javob{B}

%================= 15
\Savol{15}
$S$ ni $2\sin\frac\pi7$ ga köpaytiramiz:
\[2\sin\tfrac\pi7\,S=\left(\sin\tfrac{3\pi}7-\sin\tfrac\pi7\right)+\left(\sin\tfrac{5\pi}7-\sin\tfrac{3\pi}7\right)+\left(\sin\pi-\sin\tfrac{5\pi}7\right)=-\sin\tfrac\pi7\ \Rightarrow\ S=-\tfrac12.\]
\Javob{B}

%================= 16
\Savol{16}
$\log_2 32+27^{\log_34}=5+4^3=69$, \ $69^{\log_{69}14}=14$, \ $\frac37\cdot14=6$; \quad $4^{\log_{1/4}\frac1{10}}=4^{\log_4 10}=10$;\\
$12-2\sqrt{35}=(\sqrt7-\sqrt5)^2$ $\Rightarrow$ $-\log_{\frac19}(\sqrt7-\sqrt5)^2=\log_3(\sqrt7-\sqrt5)$, \ şuning uçun
\[\log_3\frac{81}{\sqrt7-\sqrt5}+\log_3(\sqrt7-\sqrt5)=\log_381=4.\qquad 6+10+4=20.\]
\Javob{A}

%================= 17
\Savol{17}
$f(x)=x^3-x^2-2x-1-2001$. \ $f'(x)=3x^2-2x-2=0$ $\Rightarrow$ $x=\frac{1\pm\sqrt7}3$. Lokal maksimum $x=\frac{1-\sqrt7}3\in(-1;0)$ da: $f<0$ (çunki $x^3-(x+1)^2<1$). Demak $f$ abssissa öqini faqat bir marta kesadi.
\Javob{B}

%================= 18
\Savol{18}
$-6\le|x+1|-3\le6$ $\Rightarrow$ $|x+1|\le9$ $\Rightarrow$ $-10\le x\le8$: \ $19$ ta butun son.
\Javob{C}

%================= 19
\Savol{19}
$t=\sqrt{8+x}\ge0$, $x=t^2-8$: \ $x-8t+24=(t-4)^2$, \ $16+x=t^2+8$:
\[|t-4|<\frac{t^2+8}{2(t-5)}.\]
$t<5$ da öng tomon manfiy — yeçim yöq. $t>5$ da: $2(t-4)(t-5)<t^2+8$ $\Rightarrow$ $t^2-18t+32<0$ $\Rightarrow$ $2<t<16$.\\
$5<t<16$ $\Rightarrow$ $25<8+x<256$ $\Rightarrow$ $17<x<248$: \ $248-18=230$ ta butun son.
\Javob{C}

%================= 20
\Savol{20}
$\cos2\alpha+2\sin\alpha\cos\alpha\cdot\dfrac{\sin\alpha}{\cos\alpha}=\cos2\alpha+2\sin^2\alpha=1-2\sin^2\alpha+2\sin^2\alpha=1$.
\Javob{B}

%================= 21
\Savol{21}
$f(x+1)=(x+1)^2+1$ $\Rightarrow$ $f(t)=t^2+1$. \ $f(-3)+f(1)=10+2=12$.
\Javob{B}

%================= 22
\Savol{22}
Perpendikulyarlik: $a\cdot\left(-\frac12\right)=-1$ $\Rightarrow$ $a=2$. \ $y'=\dfrac{7x/2}{\sqrt{\frac72x^2-7}}=2$ $\Rightarrow$ $\frac{49}4x^2=14x^2-28$ $\Rightarrow$ $x^2=16$, $x=4$ ($y'>0$).\\
$y(4)=\sqrt{49}=7$: urinma $y=2(x-4)+7=2x-1$. \ $a+b=2-1=1$.
\Javob{A}

%================= 23
\Savol{23}
$\sin^4u+\cos^4u=1-\frac12\sin^22u=\frac34+\frac14\cos4u$, \ $u=3x$:
\[\int_0^{\pi/12}\left(\frac34+\frac14\cos12x\right)dx=\frac34\cdot\frac\pi{12}+\frac{\sin12x}{48}\Big|_0^{\pi/12}=\frac\pi{16}.\]
\Javob{B}

%================= 24
\Savol{24}
$\overrightarrow{AE}+\overrightarrow{EC}=\overrightarrow{AC}$. \ $\overrightarrow{AC}=\overrightarrow{AB}+\overrightarrow{AD}$, \ $\overrightarrow{BD}=\overrightarrow{AD}-\overrightarrow{AB}$:
\[|\overrightarrow{AC}+\overrightarrow{BD}|=|2\overrightarrow{AD}|=16\Rightarrow AD=8,\qquad |\overrightarrow{AC}-\overrightarrow{BD}|=|2\overrightarrow{AB}|=10\Rightarrow AB=5.\]
$\triangle AED$ ning asosi $AD$, balandligi parallelogramm balandligiga teng ($E\in BC$):
\[S_{AED}=\frac12S_{ABCD}=\frac12\cdot8\cdot5\cdot\sin60^\circ=10\sqrt3.\]
\Javob{A}

%================= 25
\Savol{25}
Gipotenuzaga tuşgan mediana uning yarmiga teng: $AB=2CM=4$. Katetlardan biri $2$, ikkinçisi $\sqrt{16-4}=2\sqrt3$.
\[S=\frac12\cdot2\cdot2\sqrt3=2\sqrt3.\]
\Javob{C}

%================= 26
\Savol{26}
$E$ — kvadratning uçi, u $\angle A$ diagonalida yotadi, ya'ni $AE$ — $A$ burçak bissektrisasi:
\[\frac{AB}{AC}=\frac{BE}{EC}=\frac23,\quad AB=2k,\ AC=3k,\quad 13k^2=25^2.\]
$DE\parallel AC$: \ $\dfrac{BD}{BA}=\dfrac{DE}{AC}$ $\Rightarrow$ $\dfrac{2k-s}{2k}=\dfrac s{3k}$ $\Rightarrow$ $s=\dfrac{6k}5$.
\[s^2=\frac{36k^2}{25}=\frac{36\cdot625}{25\cdot13}=\frac{900}{13}=69\frac3{13}.\]
\Javob{B}

%================= 27
\Savol{27}
$AD=a$, $AB=b$. \ $AC^2=a^2+b^2+ab$, \ $BD^2=a^2+b^2-ab$ (kosinuslar teoremasi):
\[7(a^2+b^2+ab)=19(a^2+b^2-ab)\Rightarrow 6a^2-13ab+6b^2=0\Rightarrow \frac ab=\frac32\ \left(a>b\right).\]
\Javob{C}

%================= 28
\Savol{28}
$\sqrt{2+\sqrt3}=\frac{\sqrt6+\sqrt2}2$. Kosinuslar teoremasi:
\[BC^2=(2+\sqrt3)+3-2\cdot\frac{\sqrt6+\sqrt2}2\cdot\sqrt3\cdot\frac{\sqrt2}2=5+\sqrt3-(3+\sqrt3)=2,\quad BC=\sqrt2.\]
$R=\dfrac{BC}{2\sin45^\circ}=1$. \ $O$ dan $AC$ ga perpendikulyar vatarni teng ikkiga böladi:
\[d=\sqrt{R^2-\left(\frac{AC}2\right)^2}=\sqrt{1-\frac34}=\frac12.\]
\Javob{C}

%================= 29
\Savol{29}
Böyalgan uçburçak uçlari köpburçak uçlarida, ular orasidagi yoylar $30^\circ$, $120^\circ$, $210^\circ$; uçburçak burçaklari ularning yarmi: $15^\circ$, $60^\circ$, $105^\circ$. \ $2=2R\sin15^\circ$ $\Rightarrow$ $R=\frac1{\sin15^\circ}$.
\[S=2R^2\sin15^\circ\sin60^\circ\sin105^\circ=\frac{2\cdot\frac{\sqrt3}2\cos15^\circ}{\sin15^\circ}=\sqrt3\ctg15^\circ=\sqrt3(2+\sqrt3)=3+2\sqrt3.\]
\Javob{A}

%================= 30
\Savol{30}
$\cos\varphi=0{,}28$, $\sin\varphi=0{,}96$. \ Oğma $l=\dfrac{21}{0{,}28}=75$, perpendikulyar $h=75\cdot0{,}96=72$.
\Javob{D}

%================= 31
\Savol{31}
$y=3-x$: \ $x^2+(3-x)^2=9$ $\Rightarrow$ $2x^2-6x=0$ $\Rightarrow$ $x=0$ yoki $3$: \ $\{(3;0);\ (0;3)\}$.
\Javob{B}

%================= 32
\Savol{32}
Toq örindagi binomial koeffitsiyentlar yiğindisi barça koeffitsiyentlar yiğindisining yarmi: $\frac{2^{99}}2=2^{98}$.
\Javob{A}

%================= 33-35
\newpage
\Savol{33}
\fig{\begin{tikzpicture}[scale=1.1,font=\small]
\coordinate (O) at (0,0); \coordinate (M) at (2,0); \coordinate (S) at (0,1.1547);
\coordinate (T) at ($(O)!(O)!(S)$); \coordinate (T) at (0.5,0.866);
\draw[thick] (-2,0)--(2,0)--(S)--(-2,0); \draw[dashed] (S)--(O);
\fill[qfill] (1,0) arc (0:180:1) -- cycle; \draw (1,0) arc (0:180:1);
\draw[axs] (O)--(T); \draw[ax] (T)--(0.5,0);
\fill (O) circle (1pt); \fill (T) circle (1pt);
\node[below] at (O) {$O$}; \node[above] at (S) {$S$}; \node[below] at (M) {$M$}; \node[above right] at (T) {$T$};
\node[below] at (0.5,0) {$K$};
\draw (M) ++(150:0.45) arc (150:180:0.45); \node at ($(M)+(165:0.75)$) {\scriptsize$30^\circ$};
\node[aux,left] at (0.3,0.5) {$r$};
\end{tikzpicture}}
Piramidaning apofemasi va balandligi orqali kesim olamiz ($M$ — asos tomoni örtasi, $OM=\frac a2$). Yarim şar markazi $O$ (asos markazi), $OT\perp SM$, $OT=r$.\\
$r=OM\sin30^\circ=\frac a4$. \ $a-2r=1$ $\Rightarrow$ $a-\frac a2=1$ $\Rightarrow$ $a=2$, \ $r=\frac12$.

\textbf{33.} $r=\frac12$ \ — \textbf{A}.

\textbf{34.} $\angle TOK=90^\circ-30^\circ=60^\circ$ ($OT\perp SM$). \ $TK=OT\sin60^\circ=\frac12\cdot\frac{\sqrt3}2=\frac{\sqrt3}4$ \ — \textbf{F}.

\textbf{35.} $h=OM\tg30^\circ=\frac1{\sqrt3}$, \ $V=\frac13\cdot2^2\cdot\frac1{\sqrt3}=\frac4{3\sqrt3}=\frac{4\sqrt3}9$ \ — \textbf{E}.
\Javob{33 — A, \ 34 — F, \ 35 — E}

%================= 36
\newpage
\Savol{36}
Köpaytuvçilarga ajratamiz: $(x^2+px+q)(x^2-px+s)$, bunda $s+q-p^2=-12$, $p(s-q)=16\sqrt2$, $qs=-12$. $p=2\sqrt2$ da $s+q=-4$, $s-q=8$ $\Rightarrow$ $s=2$, $q=-6$, $qs=-12$ — töğri:
\[x^4-12x^2+16\sqrt2x-12=\left(x^2+2\sqrt2x-6\right)\left(x^2-2\sqrt2x+2\right)=\left(x^2+2\sqrt2x-6\right)\left(x-\sqrt2\right)^2.\]
$x^2+2\sqrt2x-6=0$ $\Rightarrow$ $x=-\sqrt2\pm2\sqrt2$: \ $x=\sqrt2$ yoki $x=-3\sqrt2$.

a) Turli haqiqiy ildizlar: $\sqrt2$ va $-3\sqrt2$ — \ $2$ ta.

b) $\sqrt2\cdot(-3\sqrt2)=-6$.
\Javob{a) 2; \quad b) $-6$}

%================= 37
\newpage
\Savol{37}
$\sin\left(\frac\pi2+x\right)=\cos x$. Kosinuslar yiğindisi va $1+\cos4x=2\cos^22x$:
\[\cos x+\cos3x=2\cos2x\cos x,\qquad 2\cos2x\cos x=2\cos^22x\ \Rightarrow\ \cos2x\,(\cos x-\cos2x)=0.\]
1) $\cos2x=0$ $\Rightarrow$ $x=\frac\pi4+\frac{\pi k}2$.\\
2) $\cos2x=\cos x$ $\Rightarrow$ $2x=\pm x+2\pi k$ $\Rightarrow$ $x=2\pi k$ yoki $x=\frac{2\pi k}3$, ya'ni $x=\frac{2\pi k}3$.

a) Eng kiçik musbat ildiz: $\min\left(\frac\pi4;\ \frac{2\pi}3\right)=\frac\pi4$.

b) $[0;\pi]$ da: $0$, $\frac\pi4$, $\frac{2\pi}3$, $\frac{3\pi}4$ — \ $4$ ta.
\Javob{a) $\frac\pi4$; \quad b) 4}

%================= 38
\newpage
\Savol{38}
$f(x)=(x+6)^2-6$, ya'ni $f(x)+6=(x+6)^2$. Har safar $f$ qöllanganda «$+6$» li ifoda kvadratga kötariladi:
\[f(f(x))+6=\left(f(x)+6\right)^2=(x+6)^4,\ \ldots,\ g(x)+6=(x+6)^{2^5}=(x+6)^{32}.\]
Demak $g(x)=(x+6)^{32}-6$.

a) $(x+6)^{32}=6$ $\Rightarrow$ $x+6=\pm\sqrt[32]6$ — \ $2$ ta ildiz.

b) $g(-5)=1^{32}-6=-5$.
\Javob{a) 2; \quad b) $-5$}

%================= 39
\newpage
\Savol{39}
$g(x)$ grafikka $x=p$ va $x=q$ nuqtalarda urinsa, $f(x)-g(x)=-(x-p)^2(x-q)^2$:
\[f(x)-g(x)=-\left(x^2-sx+P\right)^2=-x^4+2sx^3-(s^2+2P)x^2+2sPx-P^2,\quad s=p+q,\ P=pq.\]
Koeffitsiyentlarni tenglaymiz: $2s=-2$ $\Rightarrow$ $s=-1$; \ $-(1+2P)=3$ $\Rightarrow$ $P=-2$; \ $p,q$: $x^2+x-2=0$ $\Rightarrow$ $1$ va $-2$.\\
$x$ oldidagi: $\frac{14}3-k=2sP=4$ $\Rightarrow$ $k=\frac23$; \ ozod had: $2-l=-P^2=-4$ $\Rightarrow$ $l=6$.

a) $k+l=\frac23+6=\frac{20}3$.

b) Uriniş nuqtalari: $x=1$, $y=\frac{20}3$; \ $x=-2$, $y=\frac{14}3$.
\[d=\sqrt{3^2+2^2}=\sqrt{13}.\]
\Javob{a) $\frac{20}3$; \quad b) $\sqrt{13}$}

%================= 40
\newpage
\Savol{40}
Grafikdan: eng katta qiymat $4$ ($x=0$), eng kiçik qiymat $0$ ($x=2\pi$). \ $b+|a|=4$, $b-|a|=0$ $\Rightarrow$ $a=2$, $b=2$ ($x=0$ da maksimum). Davr $4\pi$ ($0$ dan $4\pi$ gaça bir maksimumdan keyingisigaça): $\frac{2\pi}k=4\pi$ $\Rightarrow$ $k=\frac12$.\\
$y=2\cos\frac x2+2$.

a) $\dfrac{a+b}k=\dfrac4{1/2}=8$.

b) $[2\pi;6\pi]$ — bitta töliq davr.
\[V=\pi\int_{2\pi}^{6\pi}\left(2\cos\frac x2+2\right)^2dx=4\pi\int_{2\pi}^{6\pi}\left(\cos^2\frac x2+2\cos\frac x2+1\right)dx=4\pi\left(2\pi+0+4\pi\right)=24\pi^2.\]
\Javob{a) 8; \quad b) $24\pi^2$}

%================= 41
\newpage
\Savol{41}
\fig{\begin{tikzpicture}[scale=0.42,font=\small]
\coordinate (A) at (0,0); \coordinate (D) at (10,0); \coordinate (B) at (2.8,9.6); \coordinate (C) at (12.8,9.6);
\coordinate (O1) at (2.5,1.875); \coordinate (O2) at (6.25,1.875); \coordinate (O3) at (10.3,7.725); \coordinate (O4) at (6.55,7.725);
\coordinate (Q) at (6.4,4.8);
\fill[qfill] (O1)--(O2)--(O3)--(O4)--cycle;
\draw (A)--(B)--(C)--(D)--cycle; \draw (B)--(D); \draw (A)--(C);
\foreach \o in {O1,O2,O3,O4} {\draw (\o) circle (1.875); \fill (\o) circle (2.5pt);}
\draw (O1)--(O2)--(O3)--(O4)--cycle;
\draw[ax] (O1)--(2.5,0); \draw[ax] (O2)--(6.25,0); \draw[ax] (O2)--(D);
\node[below left] at (A) {$A$}; \node[above left] at (B) {$B$}; \node[above right] at (C) {$C$}; \node[below right] at (D) {$D$};
\node[below] at (O1) {$O_1$}; \node[below] at (O2) {$O_2$}; \node[right] at (O3) {$O_3$}; \node[above] at (O4) {$O_4$};
\node[below] at (2.5,0) {$K$}; \node[below] at (6.25,0) {$L$}; \node[right] at (Q) {$Q$};
\end{tikzpicture}}
$Q$ — diagonallar kesişgan nuqta: $AQ=8$, $DQ=6$, $AD=10$. \ $\alpha=\angle QAD$: $\sin\alpha=\frac35$, $\tg\alpha=\frac34$; \ $\beta=\angle QDA$: $\tg\beta=\frac43$, \ $\tg\frac\beta2=\frac{\sin\beta}{1+\cos\beta}=\frac{4/5}{8/5}=\frac12$.\\
$O_1$ — $\angle A$ bissektrisasida (ya'ni $AC$ da), $O_2$ — $\angle ADB$ bissektrisasida. Uriniş nuqtalari $K$, $L$:
\[AK=\frac r{\tg\alpha}=\frac{4r}3,\quad KL=O_1O_2=2r,\quad LD=\frac r{\tg\frac\beta2}=2r\ \Rightarrow\ \frac{4r}3+4r=10\ \Rightarrow\ r=\frac{15}8.\]
a) $AO_1=\dfrac r{\sin\alpha}=\dfrac{5r}3=\dfrac{25}8$, simmetriya böyiça $CO_3=\dfrac{25}8$:
\[O_1O_3=AC-2\cdot\frac{25}8=16-\frac{25}4=\frac{39}4.\]
b) $O_1O_2\parallel O_4O_3$ va $O_1O_2=O_4O_3=2r$ — parallelogramm. Uning balandligi: romb balandligi $h=\frac{AC\cdot BD}{2\,AD}=\frac{16\cdot12}{20}=9{,}6$ dan $2r$ ayiriladi: $9{,}6-\frac{15}4=\frac{117}{20}$.
\[S=2r\cdot\frac{117}{20}=\frac{15}4\cdot\frac{117}{20}=\frac{351}{16}.\]
\Javob{a) $\frac{39}4$; \quad b) $\frac{351}{16}$}

%================= 42
\newpage
\Savol{42}
\fig{\begin{tikzpicture}[scale=0.45,font=\small]
\coordinate (A) at (0,0); \coordinate (B) at (7,0); \coordinate (C) at (3.5,6.062); \coordinate (P) at (2.357,1.856);
\coordinate (Q) at ($(A)+({2.357*cos(60)+1.856*sin(60)},{-2.357*sin(60)+1.856*cos(60)})$);
\draw (A)--(B)--(C)--cycle; \draw (A)--(P)--(B); \draw (P)--(C);
\draw[axs] (A)--(Q)--(B); \draw[ax] (P)--(Q);
\fill (Q) circle (1.5pt);
\node[below left] at (A) {$A$}; \node[below right] at (B) {$B$}; \node[above] at (C) {$C$};
\node[above] at (P) {$P$}; \node[below] at (Q) {$Q$};
\end{tikzpicture}}
a) $\triangle APC$ ni $A$ atrofida $60^\circ$ ga burib ($C\to B$), $P\to Q$ hosil qilamiz: $AQ=AP=3$, $BQ=CP=\sqrt{19}$, $\angle PAQ=60^\circ$ $\Rightarrow$ $\triangle APQ$ teng tomonli, $PQ=3$, $\angle APQ=60^\circ$.\\
$\triangle BPQ$: $\cos\angle BPQ=\dfrac{25+9-19}{2\cdot5\cdot3}=\dfrac12$ $\Rightarrow$ $\angle BPQ=60^\circ$, demak $\angle APB=120^\circ$:
\[AB^2=9+25+15=49,\quad AB=7,\quad P_{ABC}=21.\]
b) $\triangle APC$: $AP=3$, $PC=\sqrt{19}$, $AC=7$. \ $\cos\angle APC=\dfrac{9+19-49}{6\sqrt{19}}=-\dfrac7{2\sqrt{19}}$.\\
$P$ dan uriniş nuqtasigaça: $t=\dfrac{AP+PC-AC}2=\dfrac{\sqrt{19}-4}2$; \ $PO=\dfrac t{\cos\frac P2}$, \ $\cos^2\frac P2=\dfrac{1+\cos P}2=\dfrac{2\sqrt{19}-7}{4\sqrt{19}}$:
\[PO^2=\frac{(\sqrt{19}-4)^2}{4}\cdot\frac{4\sqrt{19}}{2\sqrt{19}-7}=\frac{(35-8\sqrt{19})\sqrt{19}}{2\sqrt{19}-7}=\frac{(35\sqrt{19}-152)(2\sqrt{19}+7)}{27}=\frac{266-59\sqrt{19}}{27}.\]
\Javob{a) 21; \quad b) $\sqrt{\frac{266-59\sqrt{19}}{27}}\approx0{,}57$}

%================= 43
\newpage
\Savol{43}
\fig{\begin{tikzpicture}[scale=1.5,font=\small]
\draw (0,0) circle (1);
\draw (90:1) \foreach \k in {1,...,11} {--({90+30*\k}:1)} --cycle;
\draw[axs] (90:1)--(150:1); \draw[axs] (90:1)--(270:1);
\node[above] at (90:1) {$A_0$}; \node[left] at (150:1) {$A_2$}; \node[below] at (270:1) {$A_n$};
\node at ($(90:1)+(-112:0.32)$) {\scriptsize$60^\circ$};
\end{tikzpicture}}
Muntazam $2n$-burçakda bir tomon $\frac{360^\circ}{2n}=\frac{180^\circ}n$ yoyni tortadi. Eng kiçik diagonal $A_0A_2$, eng kattasi — diametr $A_0A_n$. Ular orasidagi içki çizilgan burçak $A_2A_n$ yoyning yarmi:
\[\frac12\cdot(n-2)\cdot\frac{180^\circ}n=60^\circ\ \Rightarrow\ 90n-180=60n\ \Rightarrow\ n=6.\]
Köpburçak — muntazam $12$-burçak.

a) Diagonallar soni $\dfrac{12\cdot9}2=54$.

b) $S=\dfrac12\cdot12\cdot R^2\sin30^\circ=6\cdot36\cdot\dfrac12=108$.
\Javob{a) 54; \quad b) 108}

%================= 44
\newpage
\Savol{44}
\fig{\begin{tikzpicture}[scale=1.25,font=\small]
\coordinate (A) at (-0.568,3.718);
\coordinate (B) at (0.000,0.000);
\coordinate (C) at (0.984,1.609);
\coordinate (D) at (-2.687,1.178);
\coordinate (M) at (0.208,2.663);
\coordinate (N) at (-0.568,1.626);
\coordinate (K1) at (0.492,0.805);
\coordinate (K2) at (-0.852,1.393);
\coordinate (K3) at (-1.627,2.448);
\coordinate (K4) at (-0.284,1.859);
\fill[okgreen!12] (K1)--(K2)--(K3)--(K4)--cycle; \draw[okgreen] (K1)--(K2)--(K3)--(K4)--cycle;
\draw[okgreen,thick] plot[smooth cycle] coordinates {(-0.180,1.099) (-0.367,1.193) (-0.568,1.317) (-0.769,1.462) (-0.956,1.618) (-1.116,1.774) (-1.240,1.921) (-1.317,2.047) (-1.344,2.145) (-1.317,2.207) (-1.240,2.230) (-1.116,2.212) (-0.956,2.153) (-0.769,2.059) (-0.568,1.935) (-0.367,1.791) (-0.180,1.635) (-0.019,1.478) (0.104,1.332) (0.181,1.205) (0.208,1.108) (0.181,1.045) (0.104,1.022) (-0.019,1.041)};
\draw[okgreen] (M)--(0.208,1.108); \draw[okgreen] (M)--(-1.297,2.219);
\draw (A)--(B)--(C)--cycle; \draw (A)--(D)--(B); \draw[dashed] (D)--(C);
\draw[axs] (M)--(N);
\fill (M) circle (1.2pt); \fill (N) circle (1.2pt);
\node[above] at (A) {$A$}; \node[below left] at (B) {$B$}; \node[below right] at (C) {$C$}; \node[right] at (D) {$D$};
\node[right] at (M) {$M$}; \node[below] at (N) {$N$};
\end{tikzpicture}}
Kesim: $BC$, $CD$, $DA$, $AB$ qirralar örtalarini tutaştiruvçi kvadrat, tomoni $\frac a2$ (örta çiziqlar, $AC\perp BD$). Konus asosi — şu kvadratga içki çizilgan doira: $r=\frac a4$.\\
$AC$ va $BD$ ning örtalarini tutaştiruvçi $MN$ kesma ikkala qirraga ham perpendikulyar, kvadrat markazidan ötadi va kesimga perpendikulyar; $N$ — kvadrat markazi. Kesim $AC$ va $BD$ orasida örtada:
\[h=MN\cdot\frac12,\quad MN=\frac a{\sqrt2}\ \Rightarrow\ h=\frac a{2\sqrt2}.\]
Yasovçi $l=\sqrt{h^2+r^2}=\sqrt{\frac{a^2}8+\frac{a^2}{16}}=\frac{a\sqrt3}4$. \ $S_{yon}=\pi rl=\pi\cdot\frac a4\cdot\frac{a\sqrt3}4=\sqrt3\pi$ $\Rightarrow$ $a^2=16$, $a=4$.\\
$r=1$, \ $h=\sqrt2$ (bu yerda $N$ — kesim markazi, $M$ — konus uçi).

a) $V_k=\frac13\pi r^2h=\frac{\pi\sqrt2}3\approx\sqrt2$.

b) $V_t=\dfrac{a^3\sqrt2}{12}=\dfrac{64\sqrt2}{12}=\dfrac{16\sqrt2}3$.
\Javob{a) $\sqrt2$; \quad b) $\frac{16\sqrt2}3$}

%================= 45
\newpage
\Savol{45}
Tom — uçburçakli prizma, asosida $BK=8$, $KA=5$, $\angle K=60^\circ$ ($K$ — parallelepipedning yuqori old öng uçi), prizma uzunligi $12$ m.
\[BA^2=64+25-2\cdot8\cdot5\cdot\frac12=49,\quad BA=7.\]
a) Tunuka — ikki yon yoq: $(7+5)\cdot12=144$ m$^2$.

b) Parallelepiped: $8\cdot8\cdot12=768$ m$^3$ (ikki yoği $8\times8$ kvadrat). \ Tom: $\frac12\cdot8\cdot5\cdot\sin60^\circ\cdot12=10\sqrt3\cdot12=120\sqrt3$ m$^3$.
\[V=768+120\sqrt3\ \text{m}^3.\]
\Javob{a) 144; \quad b) $768+120\sqrt3$}

\ShowJavoblar
\end{document}
